\section{Beurteilung des Verlaufs des Aktienkurses}
\label{sec:kurs}

Grundlage sind die Nasdaq-Schlusskurse aus Abbildung~\ref{fig:kursverlauf}
und Tabelle~\ref{tab:kursextrema} sowie die Befunde der
Abschnitte~\ref{sec:margejeunze} und \ref{sec:reihe}.

\begin{figure}[htbp]
\centering
\begin{tikzpicture}
\begin{axis}[kursachse]
\addplot+[gray!55, thick, mark=none] table[col sep=comma, x=datum,
  y=kurs] {data/kursverlauf.csv};
\addplot+[kurspunkte] table[col sep=comma, x=datum,
  y=kurs, meta=art] {data/kursverlauf.csv};
\end{axis}
\end{tikzpicture}
\caption{Kursverlauf der Alamos-Gold-Aktie vom Jahresschluss 2020 bis zum
\KursStand{}. Jeder Punkt steht an seinem tats\"achlichen Handelstag:
{\color{kurshoch}$\blacktriangle$} Jahreshoch,
{\color{kurstief}$\blacktriangledown$} Jahrestief,
{\color{kursschluss}$\bullet$} Jahresschluss.}
\label{fig:kursverlauf}
\end{figure}

\begin{table}[htbp]
\centering
\caption{Jahresschluss-, H\"ochst- und Tiefstkurse in USD. Das Jahr 2026
reicht bis zum \KursStand{}.}
\label{tab:kursextrema}
\begin{tabular}{lrrr}
\toprule
Jahr & Schluss & H\"ochst & Tiefst\\
\midrule
\tabellenkoerper{data/kurs_tabelle}
\bottomrule
\end{tabular}
\end{table}

\FloatBarrier

\subsection{Was zu erkl\"aren ist}

Der Kurs stieg vom Jahresschluss 2021 bis zum \KursStand{} um
\Pct{\KursSeitEinundzwanzig}. Innerhalb des Jahres 2026 erreichte er am
\EinbruchTagHoch{} mit \USDje{\KursHochXXVI} sein Hoch und am
\EinbruchTagTief{} mit \USDje{\KursTiefXXVI} sein Tief. Vom Hoch bis zum Tief
verlor die Aktie \PctBetrag{\EinbruchAktie} ihres Wertes; heute steht sie
\PctBetrag{\KursAbstandHoch} unter dem Hoch und \Pct{\KursAbstandTief} \"uber
dem Tief. Gegen\"uber dem Jahresschluss 2025 liegt sie
\PctBetrag{\KursSeitFuenfundzwanzig} im Minus.

Zu erkl\"aren ist also nicht der Aufstieg, sondern der Einbruch -- und zwar
in einem Jahr, in dem der Umsatz des ersten Halbjahres um
\Pct{\HJUmsatzDelta} stieg und der freie Cash-Flow sich mehr als
verdreifachte (Abschnitt~\ref{sec:strategie}).

\subsection{Was gefallen ist: Fr\"uhindikator oder Bewertung}
\label{sec:zerlegung}

Ein Bericht, der nach einem solchen Sturz geschrieben wird, muss sagen,
\emph{was} gefallen ist. Ein R\"uckgang, der sich im Goldpreis und im
gesamten Sektor wiederfindet, ist ein Bewertungsereignis; ihn als
betriebliche Verschlechterung in die Szenarien der Flow-Analyse zu nehmen,
z\"ahlte ihn doppelt.

Tabelle~\ref{tab:zerlegung} stellt die Aktie \"uber denselben Zeitraum und
dieselben Handelstage dem Goldpreis und dem Minensektor gegen\"uber.

\begin{table}[htbp]
\centering
\caption{Zerlegung des Einbruchs 2026. Ma{\ss}st\"abe sind der SPDR Gold
Shares f\"ur den Goldpreis und der VanEck Gold Miners f\"ur den Sektor,
gerechnet \"uber dieselben Handelstage vom \EinbruchTagHoch{} bis zum
\EinbruchTagTief{}. \emph{Sekund\"arquelle} -- Kursreihen tragen keine
gedruckte Seite.}
\label{tab:zerlegung}
\begin{tabular}{lrr}
\toprule
Reihe & Hoch bis Tief & Stand heute gegen Hoch\\
\midrule
Alamos Gold & \Pct{\EinbruchAktie} & \Pct{\AbstandAktie}\\
Goldpreis & \Pct{\EinbruchGold} & \Pct{\AbstandGold}\\
Goldminen (Sektor) & \Pct{\EinbruchSektor} & \Pct{\AbstandSektor}\\
\bottomrule
\end{tabular}
\end{table}

\FloatBarrier

\bk{Die Tabelle trennt zwei Dinge, die sonst ineinanderlaufen.

\emph{Der gr\"o{\ss}ere Teil des Einbruchs war kein Alamos-Ereignis.} Der
Goldpreis fiel \"uber denselben Zeitraum um \PctBetrag{\EinbruchGold}, der
Sektor um \PctBetrag{\EinbruchSektor}. Wer einen Goldproduzenten h\"alt,
h\"alt eine gehebelte Position auf den Goldpreis; dass ein Viertel Preisverzug
gut ein Drittel Sektorverlust erzeugt, ist die Bauart dieses Gesch\"afts und
keine Nachricht \"uber das Unternehmen.

\emph{Der kleinere Teil war eines.} Alamos verlor
\PctBetrag{\EinbruchAktie}, also \Pct{\EinbruchUeberschuss} mehr als sein
eigener Sektor. Diese Differenz f\"allt zeitlich mit dem zusammen, was das
Unternehmen selbst berichtet: eine Produktion unter der revidierten Prognose,
um \Pct{\ProduktionRueckgangRund} unter dem Vorjahr, bei angehobener
Kostenprognose f\"ur 2026 (Tabelle~\ref{tab:prognose}).

\emph{Und der Abstand ist geblieben.} Seit dem Tief hat der Sektor den
Gro{\ss}teil seines Verlusts aufgeholt und steht noch
\PctBetrag{\AbstandSektor} unter dem Hoch; Alamos steht
\PctBetrag{\AbstandAktie} darunter. Die Aktie hat die Erholung ihres eigenen
Sektors also \emph{nicht} mitgemacht; der Abstand betr\"agt
\Pct{\AbstandUeberschuss}. Das ist der Teil des Kursverlaufs, den
nur betriebliche Gr\"unde erkl\"aren, und es ist der Teil, auf den ein Anleger
zu achten hat.}

\subsection{Was der Kurs nicht abbildet}

\bk{Der Kurs bildet den Anstieg der Marge je Unze ab, und er tut es mit
Hebel. Was er nach dem Befund aus Abschnitt~\ref{sec:reihe} \emph{nicht}
abbildet, ist die Frage, ob dieser Margensprung dauerhaft ist: Er stammt
\"uberwiegend aus dem Goldpreis, nicht aus sinkenden Kosten, und die Kosten
liegen im \Pct{\AISCPerzentil}-Perzentil der eigenen Zehnjahresreihe. Genau
diese Unterscheidung nimmt Teil~\ref{sec:flow} auf, indem er drei vom
Unternehmen selbst erzielte Goldpreise durchrechnet statt einen.}
